% !TEX root = ../main.tex
% PCSEL 面发射机制对比
\begin{figure}[htbp]
  \centering
  \resizebox{\textwidth}{!}{%
  \begin{tikzpicture}[x=1cm,y=1cm,>=Latex,font=\small]
    \tikzset{
      panel/.style={draw=gray!35,rounded corners=2pt,fill=gray!3},
      slab/.style={draw=gray!60,fill=gray!16},
      beam/.style={->,very thick,teal!70!black},
      guided/.style={->,very thick,blue!75!black},
      scatter/.style={->,thick,orange!90!black}
    }

    % (a) second-order Bragg extraction
    \begin{scope}[shift={(-5.2,0)}]
      \draw[panel] (-2.35,-1.55) rectangle (2.35,2.65);
      \node[font=\bfseries] at (0,2.32) {(a) 二阶 Bragg 折叠};
      \draw[slab] (-2.0,-0.35) rectangle (2.0,0.35);
      \foreach \x in {-1.55,-1.05,...,1.55} {
        \fill[white] (\x,0.08) ellipse (0.13 and 0.08);
        \draw[gray!55] (\x,0.08) ellipse (0.13 and 0.08);
      }
      \draw[guided] (-1.85,0.82)--(-0.35,0.82) node[midway,above]{导模};
      \draw[->,very thick,orange!90!black] (-0.22,0.56)--(0.48,0.56)
        node[midway,below=2pt]{$\gvec$};
      \foreach \x in {-0.42,0,0.42} {
        \draw[beam] (\x,0.35)--(\x,1.78);
      }
      \node[teal!65!black] at (0,1.98){法向辐射};
      \draw[->,thick,blue!60!black] (1.15,-0.72)--(1.85,-0.72);
      \draw[->,thick,blue!60!black] (-1.15,-0.72)--(-1.85,-0.72);
      \node[align=center,text width=4.1cm] at (0,-1.16)
      {倒格矢同时提供反馈与出光动量。};
    \end{scope}

    % (b) coherent Gamma-point band-edge emission
    \begin{scope}
      \draw[panel] (-2.35,-1.55) rectangle (2.35,2.65);
      \node[font=\bfseries] at (0,2.32) {(b) $\Gamma$ 点带边相干辐射};
      \draw[slab] (-2.0,-0.38) rectangle (2.0,0.38);
      \foreach \x in {-1.5,-1.0,...,1.5} {
        \foreach \y in {-0.15,0.15} {
          \fill[white] (\x,\y) circle (0.10);
          \draw[gray!55] (\x,\y) circle (0.10);
        }
      }
      \draw[very thick,red!75!black,domain=-1.7:1.7,samples=80]
        plot (\x,{0.58+0.16*sin(deg(3.1416*\x))});
      \foreach \x in {-1.2,-0.8,...,1.2} {
        \draw[beam] (\x,0.42)--(\x,1.72);
      }
      \draw[<->,thick,gray!75] (-0.95,1.30)--(-0.45,1.30);
      \node[gray!60,fill=white,inner sep=1pt] at (-0.70,1.53){同相};
      \node[align=center,text width=4.0cm] at (0,-1.16)
      {带边慢光增强驻波场，远场相位更有序。};
    \end{scope}

    % (c) defect or disorder assisted leakage
    \begin{scope}[shift={(5.2,0)}]
      \draw[panel] (-2.35,-1.55) rectangle (2.35,2.65);
      \node[font=\bfseries] at (0,2.32) {(c) 缺陷/无序辅助散射};
      \draw[slab] (-2.0,-0.35) rectangle (2.0,0.35);
      \foreach \x in {-1.55,-1.05,-0.55,0.85,1.35} {
        \fill[white] (\x,0.08) ellipse (0.13 and 0.08);
        \draw[gray!55] (\x,0.08) ellipse (0.13 and 0.08);
      }
      \fill[red!22] (0.05,-0.28) ellipse (0.32 and 0.20);
      \draw[red!75!black,densely dashed,thick] (0.05,-0.28) ellipse (0.32 and 0.20);
      \node[red!75!black,fill=white,inner sep=1pt] at (0.28,0.52){局域扰动};
      \draw[guided] (-1.85,0.82)--(-0.15,0.82);
      \foreach \ang in {35,65,100,135,160} {
        \draw[scatter] (0.05,0.02)--++(\ang:1.15);
      }
      \node[orange!80!black,align=center,font=\scriptsize,fill=white,inner sep=1pt]
        at (0,1.62){分散出光\\相干性较弱};
      \node[align=center,text width=4.1cm] at (0,-1.16)
      {局域散射可增加出光，但常伴随损耗与远场劣化。};
    \end{scope}
  \end{tikzpicture}%
  }
  \caption{光子晶体面发射激光器中三类主要面发射机制对比：二级衍射、带边相干辐射与缺陷辅助散射。三者都可以理解为周期或局域扰动把面内导模耦合到辐射通道，但相干性、方向性和损耗代价不同。}
  \label{fig:ch06_surface_emission_mechanisms}
\end{figure}
